
\begin{corollaire}
\label{XIII.1.9}
Sean $f\colon X\to Y$ un $S$-morfismo entero y $g\colon Y\to Z$ un
$S$-morfismo. Si~$F$ es un haz de conjuntos sobre~$X$, entonces
$(f_*F,g)$ es cohomológicamente propio relativamente a~$S$ en dimensión
$\leqslant -1$ (\resp $\leqslant 0$) si y solo si lo mismo sucede con
$(F,gf)$. Si~$F$ es un haz de grupos sobre~$X$, entonces $(f_*F,g)$ es
cohomológicamente propio relativamente a~$S$ en dimensión
$\leqslant -1$ (\resp $\leqslant 0$, \resp $\leqslant 1$) si y solo si
lo mismo sucede con~$(F,gf)$.
\end{corollaire}

La afirmación
\marginpar{357}
relativa al caso de un haz de conjuntos se deduce de~\Ref{XIII.1.6} y
del hecho de que $(F,f)$ es cohomológicamente propio relativamente a~$S$
en dimensión $\leqslant 0$. Sean~$F$ un haz de grupos sobre~$X$ y~$\Phi$
el stack de los torsores bajo~$F$. Según SGA~4 VIII~5.8, todo torsor
bajo~$F$ es localmente trivial sobre~$Y$. De ello se sigue que el stack
$f_*\Phi$ es equivalente al stack de los torsores bajo~$f_*F$; la
equivalencia se obtiene asociando a todo esquema~$Y_1$ étale sobre~$Y$ y
a todo torsor~$P$ sobre~$X_1=X\times_Y Y_1$ bajo~$F|X_1$ el torsor
$f_*P$ bajo~$f_*F|Y_1$. Como $(F,f)$ es cohomológicamente propio
relativamente a~$S$ en dimensión $\leqslant 1$, el corolario se deduce
de~\Ref{XIII.1.6}.
